\section{Scope and further questions}
The global five-term classification, the infinite exact family and the density benchmark use classical complex multiplication and local analysis. The uniform near-degree lower bound uses Dirichlet GRH for small inert primes, class numbers and coefficient heights. The Frobenius optimization uses the one-sided residue argument of the Lean proof. Its application to a nonabelian splitting field leaves an analytic inequality open. The archive includes source identities, finite witnesses, replay code and the recorded output of the checks.

\begin{enumerate}
\item Does every singular modulus of degree $h\ge6$ require at least \emph{seven} terms? This is the next support threshold after Theorem~\ref{thm:four}.
\item Can one bound $s(\Delta)$ below by a fixed positive fraction of $h(\Delta)$ for every order? The current argument loses a factor of order $p$ when the conjugates occupy different supersingular residues.
\item Can one determine the true relative density of Gaussian prime conductors with $s=h+1$? The value in Section~\ref{sec:family} counts a sufficient factorial inequality. Determinants can remain nonzero even when two coefficients have the same valuation.
\item For the dihedral field in Section~\ref{sec:transfer}, can a comparison of arithmetic quantities give the required one-sided quartic-nonresidue estimate? The normalized character coefficient is already optimal within the class of tests analyzed here.
\end{enumerate}

\appendix
\section{Replay of the earlier four-term reduction}\label{sec:earlier}
The argument below records the earlier route in \cite{prior} to the 79 discriminants treated in Section~\ref{sec:finite}. The global five-term proof in Sections~\ref{sec:global-five} and \ref{sec:four-sieves} works independently of this analytic cutoff.
\paragraph{Four different sizes.}
Suppose four nonzero conjugates have increasing absolute values $r_i$ with $r_{i+1}/r_i\ge16$. For exponents $e_1<\cdots<e_4$, multiply column $j$ of $(z_i^{e_j})$ by a positive $c_j$ with $c_{j+1}/c_j=(r_jr_{j+1})^{-(e_{j+1}-e_j)/2}$. On either side of the diagonal the absolute entries decrease by factors of at least four as one moves away from it. The sum of the other entries in each row is at most $2/3$ of the diagonal entry in absolute value. Thus the matrix is nonsingular. Any relation with fewer terms could be padded by unused exponents. For tighter gaps the earlier checker instead sums all 23 nondiagonal permutation magnitudes relative to the diagonal. The sum is largest at unit exponent gaps, by rearrangement, and a bound strictly below one suffices.

\paragraph{The analytic cutoff.}
Put $D=|\Delta|$ and $R=\sqrt D$. A split prime $p\nmid D$, $p\le R/2$, gives a primitive reduced form with leading coefficient $p$. Indeed one chooses $b^2\equiv\Delta\pmod{4p}$, $|b|\le p$, and puts $c=(b^2+D)/(4p)\ge p$. The absolute value of its singular modulus is within $2079$ of $\exp(\pi R/p)$ by the cusp estimate in \cite{riffaut}. The principal form gives leading coefficient 1. For $D\ge10^{11}$ three eligible primes occur, respectively, in $(0,.046R]$, $(.048R,.189R]$, $(.230R,.500R]$. To record the constants, set
\[
 \begin{split}
 E(y)&=\sqrt y\big[(\log y/(2\pi)+2)\log D+2(\log^2 y/(8\pi)+2)\big],\\
 B(y)&=2\sqrt y\log y+\log D+2\log 2,\qquad E(0)=0 .
 \end{split}
\]
Greni\'e and Molteni's prime ideal estimate for $K$ gives $|\psi_K(y)-y|\le E(y)$ for $y\ge100$. Its hypothesis follows from Dirichlet GRH. Over rational primes $p\le\sqrt y$, prime ideal powers contribute at most $2\sqrt y\log y$. The ramified primes and the split primes dividing $f$ cost at most $\log|d_K|+2\log f=\log D$ at larger primes. The extra term allows removal of $2$. Thus twice the eligible prime log-weight on $(x,y]$ is at least $y-x-E(x)-E(y)-B(y)$. At the three intervals the certified lower margins are respectively $68,191,751$ at $D=10^{11}$. After division by $R$ the error terms decrease: each monomial in $E(cR)/R$ contains $R^{-1/2}$ and at most two logarithm factors, each at least 9 for the positive endpoints. The same monotonicity holds for $B(cR)/R$. For two consecutive prime intervals with adjacent endpoints $vR,u'R$, the cusp ratio is bounded below by $(e^{\pi/v}-2079)/(e^{\pi/u'}+2079)>16$; the principal-to-first ratio is even larger since $p\ge3$. This proves the large range.\cite{gm,prior}

\paragraph{The finite range.}
For $10^7\le D<10^{11}$, enumerate entire residue classes where $1,p,p^2,p^3$ give four eligible reduced forms for a split prime $p$ with $4p^6\le D$. Include the powers of 2 when $D\equiv7\pmod 8$. Lifting solutions of $b^2\equiv\Delta\pmod{4p^j}$ produces the forms. For consecutive leading coefficients $a<b$, the tests $D\ge16a^2$ and $9D(b-a)^2\ge16a^2b^2$ imply a ratio exceeding 16 by the cusp bound. They hold for these powers in the specified range. The rest are tested individually. The CRT sieve accounts for all $49{,}995{,}000{,}000$ discriminants in the half-open range, including $47{,}116{,}196$ settled with three individually tested split primes. Each bin balances against the exact count in residue classes 0 and 3 modulo 4.

Below $10^7$, the integer form sieve and interval determinant test leave 80 inputs of degree at least four. The reduction enumerates primitive forms with $0\le b\le a\le c$ for absolute values, treating nonreal conjugate pairs once. The published output also covers the endpoint $10^7$. There are 67 low-degree inputs, $4{,}999{,}603$ exclusions at the first pass and another 250 by rigorous intervals. At $\Delta=-4075$ a direct complex determinant removes the sole remaining sextic. In degree six there are three smaller conjugates with moduli in $(2365,2366)$, $(2846,2847)$ and $(3701,3703)$, together with a principal conjugate larger than $10^{80}$ in modulus. If either of the first two exponent gaps is at least 12, diagonal dominance of the determinant expansion suffices. For first gaps 1 through 11, dividing by the principal root's highest power reduces to 121 nonzero complex minors, with a uniformly smaller error for every last gap. The small sieve, large sieve, boundary numbers, and all 121 minors were rerun. The archive includes the original checkers and expected outputs.\cite{prior}


\begin{thebibliography}{99}\small\raggedright
\bibitem{repo} L. Sernau and collaborators, \emph{GRH for Dirichlet $L$-functions, in Lean 4}. \href{https://github.com/lsernau/grh-lean/tree/72fead102f663d5c7a407964a13dc0cf9d977207}{Source revision}, \href{https://github.com/lsernau/grh-lean/blob/72fead102f663d5c7a407964a13dc0cf9d977207/grh/src/Statements.lean}{public types}, \href{https://github.com/lsernau/grh-lean/blob/72fead102f663d5c7a407964a13dc0cf9d977207/grh/src/New/Reduction/Zeros/Criterion.lean}{one-sided criterion}, \href{https://github.com/lsernau/grh-lean/blob/72fead102f663d5c7a407964a13dc0cf9d977207/grh/GUIDE.md}{reading guide}.
\bibitem{mitcheck} A. Sutherland, \emph{Verification files for the shorter Lean proof of RH and GRH}, reproductions and comparison logs for both editions. \url{https://math.mit.edu/~drew/shortgrh/}.
\bibitem{prior} V. Reshetnikov, ProveIt, \href{https://github.com/VladimirReshetnikov/ProveIt/blob/c477d08ef7be4241e94bd1cffb43ff56240321e3/docs/incoming/grh_cm_prime_applications.zip}{prior research archive}, especially \texttt{riffaut.tex}, \texttt{cm-finite.tex}, \texttt{cm\_small.tex}, \texttt{cm4075.tex}, \texttt{cm-conclusion.tex}. Source identifiers and a comparison of earlier archives are recorded in \texttt{checks/prior\_audit.md}.
\bibitem{companion} Companion report, \emph{Sparse equations for singular moduli and the sign of weighted prime sums}, 25 November 2026, \texttt{grh\_cm\_four\_terms\_and\_prime\_bias.zip}. See its quintic section and further research questions 1 and 2. The file identity is in \texttt{checks/prior\_audit.md}.
\bibitem{sparseparallel} Parallel report, \emph{Sparse equations for singular moduli}, 25 November 2026, \texttt{grh\_cm\_sparse\_equations.zip}. See its five-term reduction, verification through width $65{,}536$, and extremal family. File identity appears in \texttt{checks/prior\_audit.md}.
\bibitem{riffaut} Y. Bilu, F. Luca and A. Pizarro-Madariaga, \emph{Trinomials, singular moduli and Riffaut's conjecture}, Annali Sc. Norm. Sup. Pisa (5) 23 (2022), 2003--2048. \href{https://arxiv.org/abs/2003.06547}{arXiv:2003.06547}.
\bibitem{aven} M. Avenda\~no and T. Krick, \emph{Sharp bounds for the number of roots of univariate fewnomials}, J. Number Theory 131 (2011), 1209--1228, esp. Prop. 3.8 and the discussion of Lenstra's local bound. \href{https://arxiv.org/abs/1008.4808}{arXiv:1008.4808}.
\bibitem{cmred} J. Belding, R. Br\"oker, A. Enge and K. Lauter, \emph{Computing Hilbert class polynomials}, esp. \S3 on inert primes and CM reduction. \href{https://arxiv.org/abs/0802.0979}{arXiv:0802.0979}. See also D. Cox, \emph{Primes of the Form $x^2+ny^2$}, 2nd ed., on ring class fields, class numbers and conjugates.
\bibitem{supersing} J. Silverman, \emph{The Arithmetic of Elliptic Curves}, 2nd ed., Springer (2009), Chapter V, Theorem 4.1, supersingular invariants.
\bibitem{suth} A. V. Sutherland, \emph{Computing Hilbert class polynomials with the Chinese remainder theorem}, Math. Comp. 80 (2011), 501--538, Appendix 1, Lemma 8. \href{https://arxiv.org/abs/0903.2785}{arXiv:0903.2785}.
\bibitem{watkins} M. Watkins, \emph{Class numbers of imaginary quadratic fields}, Math. Comp. 73 (2004), 907--938, Table 4. \href{https://doi.org/10.1090/S0025-5718-03-01517-5}{Publisher's version}.
\bibitem{gm} L. Greni\'e and G. Molteni, \emph{Explicit versions of the prime ideal theorem for Dedekind zeta functions under GRH}, Math. Comp. 85 (2016), 889--906, Corollary 1.2. \href{https://arxiv.org/abs/1312.4463}{arXiv:1312.4463}.
\bibitem{gms} L. Greni\'e and G. Molteni, \emph{Explicit smoothed prime ideals theorems under GRH}, Theorem 1.1. \href{https://arxiv.org/abs/1312.4465}{arXiv:1312.4465}.
\bibitem{pollack} P. Pollack, \emph{Bounds for the first several prime character nonresidues}, Theorem 1.1. \href{https://arxiv.org/abs/1508.05035}{arXiv:1508.05035}.
\end{thebibliography}
